\protect \section *{\nameref *{sec:fusion}}
\begin{Solution}{10.{1}}
		$P_f\approx0{,}80$, $P_\alpha\approx0{,}16$, $P_{br}\approx0{,}017$~Вт/см$^3$; $\tau_E\approx3$~с ($n\tau_E\approx3\times10^{14}~\mathrm{см}^{-3}\mathrm{с}$); при 20~кеВ $n\tau_E\approx1{,}6\times10^{14}$.
	
\end{Solution}
\begin{Solution}{10.{2}}
		$v_\alpha\approx1{,}3\times10^{9}$~см/с, $r_{L\alpha}\approx5{,}1$~см; частка $Q/(Q+5)$, при $Q=10$ це $2/3$.
	
\end{Solution}
\begin{Solution}{10.{3}}
		$Q_{\min}\approx5{,}3$; з~урахуванням власних потреб умова стає
	$0{,}8\,\eta_{\mathrm{th}}[(0{,}8M+0{,}2)Q+1]\geq1/\eta_{\mathrm{aux}}$, звідки $Q_{\min}\approx6{,}8$.
	
\end{Solution}
\begin{Solution}{10.{4}}
		$R=10$~мкм, $c_s\sim9\times10^{7}$~см/с, $\tau\sim10$~пс.
	
\end{Solution}
\begin{Solution}{10.{5}}
		D--T дає $17{,}6$~МеВ, із~них лише $20\%$ у~заряджених продуктах;
	$\mathrm D+{}^3\mathrm{He}\to{}^4\mathrm{He}\,(3{,}6~\text{МеВ})+p\,(14{,}7~\text{МеВ})$ дає $18{,}3$~МеВ,
	усе в~заряджених частинках. Оптимальна температура для D--$^3$He значно вища
	($\sim50$--$100$~кеВ проти $10$--$20$~кеВ), а~$\langle\sigma v\rangle$ там приблизно на~порядок менше
	за~D--T при $10$--$20$~кеВ. Перевага: мало нейтронів (лише від побічних D--D-реакцій), отже менша
	активація конструкцій і~можливість прямого перетворення енергії заряджених продуктів в~електрику.
	Недоліки: потрібні вищі $T$ і~$nT\tau_E$, гальмівне випромінювання ($Z=2$) відносно більше,
	а~$^3$He на~Землі майже немає.
	
\end{Solution}
\begin{Solution}{10.{6}}
		$3{,}4\times10^{14}$~Дж/кг; якщо все тепло $P_{\mathrm{th}}=2{,}5$~ГВт дає синтез,~---
	$\approx0{,}64$~кг/добу (з~урахуванням тепла бланкета, $M=1{,}2$,~--- $\approx0{,}55$~кг/добу) проти
	$\approx8{,}5$~тис.~т вугілля на~добу.
	
\end{Solution}
\begin{Solution}{10.{7}}
		$dN_n/dt=P_f/(17{,}6~\text{МеВ})\approx1{,}8\times10^{20}$~с$^{-1}$; $\Delta E\approx82{,}5\sqrt{15}\approx320$~кеВ.
	
\end{Solution}
